\documentclass[11pt]{article}
\usepackage[utf8]{inputenc}
\usepackage[T1]{fontenc}
\usepackage{amsmath}
\usepackage{amssymb}
\usepackage{amsthm}
\usepackage[margin=1.15in]{geometry}
\usepackage{booktabs}
\usepackage{hyperref}

\newtheorem{theorem}{Theorem}
\newtheorem{lemma}{Lemma}
\theoremstyle{remark}
\newtheorem{remark}{Remark}

\title{Refutation of conjecture 00000000471\\[2mm]
\large The hypertree count of $K_n^{(3)}$: a closed form that exceeds the
number of subhypergraphs}
\author{Submitted to The Last Math Competition}
\date{\today}

\begin{document}
\maketitle

\begin{abstract}
\noindent
Conjecture 00000000471 proposes the closed form
$t(K_n^{(3)}) = n^{\binom{n-1}{2}-1}\prod_{i=1}^{n-1}(i^2-i+1)$ for the number of
hypertrees of the complete $3$-uniform hypergraph, and states that it covers all
known values for $n \le 6$. We give two independent refutations. First, a
hypertree is a subhypergraph and is therefore determined by a subset of the
$\binom{n}{3}$ edges, so $t(K_n^{(3)}) \le 2^{\binom{n}{3}}$; the conjectured
formula violates this bound at every $n$ with $3 \le n \le 13$, and in particular
at all six values the conjecture claims to cover. This argument uses nothing
about connectedness or minimality. Second, exhaustive enumeration gives
$t(K_3^{(3)}) = 1$, $t(K_4^{(3)}) = 6$ and $t(K_5^{(3)}) = 25$, against $3$,
$336$ and $853125$ from the formula. Both refutations are formalised in core
Lean~4.
\end{abstract}

\section{The conjecture}

The statement under consideration, verbatim from the corpus, is the following.

\begin{quote}
\emph{Definition: A hypertree of the $r$-uniform complete hypergraph
$K_n^{(r)}$ is a minimally connected subhypergraph of its edge set;
$t(K_n^{(r)})$ denotes their count. Conjecture: For $r = 3$,
$t(K_n^{(3)}) = n^{\binom{n-1}{2}-1}\cdot\prod_{i=1}^{n-1}f(i)$ with the
explicit polynomial $f(i) = i^2 - i + 1$; this formula covers all known values
for $n \le 6$.}
\end{quote}

Write $V = \{1,\dots,n\}$ and let $E_n$ be the set of all $\binom{n}{3}$
three-element subsets of~$V$, so that $K_n^{(3)} = (V, E_n)$. A
\emph{subhypergraph} of $K_n^{(3)}$, in the only sense consistent with the
phrase ``subhypergraph of its edge set'', is a pair $(V, F)$ with $F \subseteq
E_n$; it is determined by~$F$. A subhypergraph is \emph{connected} if its edges
cover $V$ and the graph obtained by replacing each edge by a triangle on its
three vertices is connected; it is \emph{minimally connected} if it is
connected and deleting any single edge destroys connectedness. A \emph{hypertree}
is a minimally connected subhypergraph, and $t(K_n^{(3)})$ counts them.

We write
\[
  \Phi(n) \;=\; n^{\binom{n-1}{2}-1}\prod_{i=1}^{n-1}(i^2-i+1)
\]
for the conjectured value.

\section{Refutation I: the counting bound}

This is the load-bearing argument, and it is degenerate in its simplicity.

\begin{lemma}\label{lem:bound}
For every $n \ge 3$, $\;t(K_n^{(3)}) \le 2^{\binom{n}{3}}$.
\end{lemma}

\begin{proof}
A hypertree is a subhypergraph of $K_n^{(3)}$, hence is determined by a subset
$F \subseteq E_n$ of its edges. There are $2^{|E_n|} = 2^{\binom{n}{3}}$ such
subsets. The map $F \mapsto (V,F)$ is injective on subhypergraphs, so there are
at most $2^{\binom{n}{3}}$ of them, and \emph{a fortiori} at most that many
hypertrees.
\end{proof}

Lemma~\ref{lem:bound} does not mention connectedness, minimality, $r = 3$ beyond
the edge count, or anything about the structure of hypergraphs. It survives every
possible reading of the two undefined-looking words in the conjecture.

\begin{theorem}\label{thm:bound}
$\Phi(n) > 2^{\binom{n}{3}}$ for every $n$ with $3 \le n \le 13$. In particular
the conjecture fails at $n = 3$.
\end{theorem}

\begin{proof}
For $n = 3$ we have $\binom{n-1}{2} - 1 = \binom{2}{2} - 1 = 0$ and
$\prod_{i=1}^{2}(i^2-i+1) = 1 \cdot 3 = 3$, so $\Phi(3) = 3^0 \cdot 3 = 3$,
while $2^{\binom{3}{3}} = 2^1 = 2$. Thus $\Phi(3) = 3 > 2$, contradicting
Lemma~\ref{lem:bound}.

The remaining values are tabulated in Table~\ref{tab:bound}; they are finite
computations with integers. The bound is violated for $3 \le n \le 13$ and holds
from $n = 14$ onwards, where $\binom{n}{3} \approx n^3/6$ finally outgrows
$\bigl(\binom{n-1}{2}-1\bigr)\log_2 n + \sum_{i<n}\log_2(i^2-i+1) \approx
\frac{n^2}{2}\log_2 n$.
\end{proof}

\begin{table}[h]
\centering
\begin{tabular}{rrrrr}
\toprule
$n$ & $\binom{n}{3}$ & $2^{\binom{n}{3}}$ & $\Phi(n)$ & $\Phi(n)/2^{\binom{n}{3}}$ \\
\midrule
3 & 1   & 2        & 3            & $1.5\times$ \\
4 & 4   & 16       & 336          & $21\times$ \\
5 & 10  & 1024     & 853125       & $833\times$ \\
6 & 20  & 1048576  & 57775431168  & $55099\times$ \\
\midrule
11 & 165 & $\approx 5.0\cdot 10^{49}$  & $\approx 2.0\cdot 10^{58}$  & $\approx 4\cdot 10^{8}\times$ \\
12 & 220 & $\approx 1.6\cdot 10^{66}$  & $\approx 6.3\cdot 10^{72}$  & $\approx 4\cdot 10^{6}\times$ \\
13 & 286 & $\approx 1.3\cdot 10^{86}$  & $\approx 7.9\cdot 10^{88}$  & $\approx 8\cdot 10^{2}\times$ \\
14 & 364 & $\approx 4.0\cdot 10^{109}$ & $\approx 1.3\cdot 10^{107}$ & bound holds \\
\bottomrule
\end{tabular}
\caption{The conjectured formula against the number of subhypergraphs. The
conjecture asserts agreement for $n \le 6$; every one of those values is in the
left-hand part of the table.}
\label{tab:bound}
\end{table}

\begin{remark}
One counterexample suffices to refute a universal claim, and $n = 3$ is the
smallest case in which $K_n^{(3)}$ is nonempty at all. It is worth noting how
far the failure extends: the formula is not merely wrong at one value, it is
impossible throughout the range the conjecture explicitly claims to have checked.
\end{remark}

\section{Refutation II: exact enumeration}

The counting bound shows the formula is too large. Enumerating directly shows
what the true values are.

\paragraph{$n = 3$.} Here $|E_3| = \binom{3}{3} = 1$; the unique edge is
$\{1,2,3\}$. There are two subhypergraphs, the empty one and the whole one. The
empty one is not connected (its edges do not cover $V$), and the whole one is
connected and vacuously minimal, since deleting the only edge disconnects it.
Hence
\[
  t(K_3^{(3)}) = 1, \qquad \Phi(3) = 3.
\]

\paragraph{$n = 4$.} Here $|E_4| = 4$. A single edge covers only three of the
four vertices, so a connected subhypergraph needs at least two edges; two
distinct triples cover all four vertices exactly when they omit different
vertices, giving $\binom{4}{2} = 6$ pairs, each of which is connected and
minimal. No set of three or more edges is minimal. Hence
\[
  t(K_4^{(3)}) = 6, \qquad \Phi(4) = 4^2 \cdot (1\cdot 3\cdot 7) = 336.
\]

\paragraph{$n = 5$.} Here $|E_5| = 10$ and there are $2^{10} = 1024$
subhypergraphs. Exhaustive search gives
\[
  t(K_5^{(3)}) = 25, \qquad
  \Phi(5) = 5^5 \cdot (1\cdot 3\cdot 7\cdot 13) = 853125.
\]
The $25$ split as $15 + 10$: fifteen consist of two edges meeting in a single
vertex and covering all five vertices, and ten consist of three edges sharing a
common pair, the third entries running over the remaining three vertices.

\begin{table}[h]
\centering
\begin{tabular}{rrrr}
\toprule
$n$ & subhypergraphs & $t(K_n^{(3)})$ & $\Phi(n)$ \\
\midrule
3 & 2    & 1  & 3       \\
4 & 16   & 6  & 336     \\
5 & 1024 & 25 & 853125  \\
\bottomrule
\end{tabular}
\caption{Exact counts against the conjectured formula.}
\label{tab:exact}
\end{table}

The value $n=6$ has $2^{20} = 1048576$ subhypergraphs; we do not enumerate it,
because Theorem~\ref{thm:bound} already disposes of it.

\section{Why no re-reading of the definition rescues it}

The words ``connected'' and ``minimal'' admit some variation, so it is worth
isolating what each refutation depends on.

\begin{itemize}
\item Theorem~\ref{thm:bound} uses only that a hypertree is determined by a
subset of the edge set. Any definition of ``minimally connected
subhypergraph'' for which this is true --- which is to say, any definition at
all --- is refuted. In particular it does not matter whether connectivity is
read via the $2$-section, via incidence-graph connectivity, or via any stronger
notion, and it does not matter whether ``minimal'' means minimal by inclusion or
of minimum cardinality.
\item The enumeration at $n = 3$ uses only that the empty subhypergraph is not
connected and that the full one is. It is insensitive to the same variations.
\end{itemize}

The only reading we know of that would evade Lemma~\ref{lem:bound} is one in
which a subhypergraph may also delete vertices, so that $t$ counts hypertrees
spanning \emph{any} subset of $V$ rather than all of $V$. Under that reading the
notation $t(K_n^{(r)})$ would be misleading (it would depend only on $n$ through
the choice of available vertices), and the count at $n = 3$ would be $1$ for the
spanning hypertree plus one for each singleton vertex, i.e.\ $4$, which still
disagrees with $\Phi(3) = 3$. We record this only for completeness; the
standard reading, and the one consistent with the classical notation
$t(K_n) = n^{n-2}$ for graphs, is the spanning one.

\section{Verification}

\begin{itemize}
\item \texttt{python3 reproduce.py} --- standalone, standard library only.
Recomputes Table~\ref{tab:bound} and Table~\ref{tab:exact} from scratch by
enumerating all $2^{\binom{n}{3}}$ subhypergraphs for $n = 3,4,5$ and checking
connectedness and minimality of each.
\item \texttt{lean4/} --- a core Lean~4 project (no Mathlib, no \texttt{sorry}).
Run \texttt{lake build \&\& lake env lean Check.lean}.
\end{itemize}

The formalisation contains
\begin{itemize}
\item \texttt{sublists\_length}: the powerset of an $m$-element set has $2^m$
elements;
\item \texttt{filter\_length\_le}: filtering does not lengthen a list, hence
\texttt{hypertrees\_le\_subhypergraphs}: $t(K_n^{(3)}) \le 2^{\binom{n}{3}}$;
\item \texttt{refutation\_bound\_3}: $\Phi(3) > t(K_3^{(3)})$, obtained by
combining the bound with the exact value $2^{\binom{3}{3}} = 2$;
\item \texttt{refutation\_bound\_6}: $\Phi(6) > 2^{\binom{6}{3}}$, i.e.\
$57775431168 > 2^{20}$;
\item \texttt{hypertrees3\_length}, \texttt{hypertrees4\_length},
\texttt{hypertrees5\_length}: the exact counts $1$, $6$, $25$, each proved by
\texttt{rfl}, so that the kernel itself performs the exhaustive search;
\item \texttt{refutation\_exact\_3/4/5}: $\Phi(n) \ne t(K_n^{(3)})$ for
$n = 3,4,5$.
\end{itemize}

Axiom audit (\texttt{\#print axioms}): \texttt{hypertrees3\_length},
\texttt{hypertrees4\_length}, \texttt{hypertrees5\_length} and the three
\texttt{refutation\_exact\_*} theorems \emph{depend on no axioms at all}. The
bound theorems depend on \texttt{propext} and \texttt{Quot.sound}, inherited
from the standard treatment of lists.

\section{Scope}

We refute the closed form as stated. We make no claim about the correct value of
$t(K_n^{(3)})$ beyond $n \le 5$, and no claim about the analogous formula for
$r \ne 3$. The three exact values $1, 6, 25$ are reported because they pin down
the small cases, not because they suggest a pattern; three terms are not enough
to guess one.

\end{document}
